\documentclass[10pt,a4paper]{amsart}
\usepackage[margin=19mm]{geometry}
\usepackage{amsmath,amssymb,mathtools}
\usepackage[scr=rsfs,cal=euler]{mathalfa}
\usepackage{xfrac}
\usepackage[unicode,hidelinks,bookmarks=false]{hyperref}
\newcommand{\ie}{i.e.\ }
\newcommand{\cf}{cf.\ }
\newcommand{\resp}{respectively\ }
\newcommand{\Subsection}{\subsection}
\providecommand{\nobreakdash}{\nobreak\hbox{-}}
\setlength{\emergencystretch}{2em}
\multlinegap0pt
\usepackage{amssymb}
\newtheorem{corollary}{Corollary}
\newtheorem{theorem}{Theorem}
\newtheorem{lemma}{Lemma}
\def\diag{{\rm diag}}
\def\rddots{\cdot^{\displaystyle\cdot^{\displaystyle\cdot}}}
\title{Numerical radius of certain two-by-two block matrices}
\author{Hwa-Long Gau, Jia-Huo Hong, Chi-Kwong Li, Kuo-Zhong Wang}
\date{Source: arXiv:2606.08576v1 (CC BY 4.0)}
\begin{document}
\maketitle

\noindent\emph{Source and licence.} Numerical radius of certain two-by-two block matrices. Version arXiv:2606.08576v1 (CC BY 4.0), distributed by arXiv under the Creative Commons Attribution 4.0 International licence, http://creativecommons.org/licenses/by/4.0/, as recorded in the arXiv licence metadata for this exact version and shown on the abstract page https://arxiv.org/abs/2606.08576v1. Full licence notice: \texttt{licenses/arxiv-2606.08576-CC-BY-4.0.txt}.\par\smallskip

\noindent\textit{Editorial scope.} Counted unit: the article's theorem t36 (source0.tex lines 817--831). The article's complete original proof of it is delivered with it (source0.tex lines 892--1006, one \texttt{proof} environment of 115 lines, which the article places away from the statement). No mathematics is written, repaired or re-derived: the statement and the proof are the article's own lines, in the article's own notation, and no retained line is substituted or re-flowed.  Retained uncounted as the source-local context the proof needs: lines 270--275; lines 810--812; lines 835--843.  Nothing was omitted from any retained span, so the package carries no omission record.  No retained line contains a citation, so the wrapper carries no bibliography and no bibliography entry.  No retained line contains a figure environment, an image-inclusion command or drawing code, so no image is delivered and the package declares no asset dependency.  Editorial formatting only, in addition: an AMS \texttt{amsart} preamble that restates the article's own declarations and macros used by the retained lines, namely the environments \texttt{corollary}, \texttt{theorem}, \texttt{lemma} on separate counters, together with the standard packages those lines need.  The retained environments print in this document as Corollary 1, Theorem 1, Lemma 1 in order of appearance, whereas the article prints the same items with the numbering of its own full document.  The complete native source of the article as distributed in the arXiv e-print for this exact version is bundled unchanged as \texttt{sources/202340/source0.tex}.
\begin{corollary} \label{2.4} Let
$T = \begin{pmatrix} A & B \cr 0 & 0 \cr\end{pmatrix} \in B(H \oplus H)$.
If $B_1, B_2\in B(H)$ such that $B_1B_1^* \le BB^* \le B_2B_2^*$ and $T_j$ is obtained
from $T$ be replacing $B$ with $B_j$ for $j = 1,2$, then
$$W(T_1) \subseteq W(T) \subseteq W(T_2).$$
\end{corollary}

\begin{equation}\label{eq1}
S_B=\begin{pmatrix} s_1 & & \\ & \ddots & \\ & & s_k\end{pmatrix}\in M_k \ \ \ \mbox{and} \ \ \ J_k=\begin{pmatrix} & & 1\\ & \rddots & \\ 1 & & \end{pmatrix}\in M_k.
\end{equation}

\begin{lemma}\label{l35}
Let $1\le q\le p$, $k=p+q$, $B=\diag(s_1, \dots, s_k)\in M_k$ with $s_1\ge s_2\ge \cdots\ge s_k\ge 0$ and $B_1=\diag(s_{2}, \dots, s_{k-1})\in M_{k-2}$. Then, for any unitary $U\in M_k$, we have
$$
 w\Big(\begin{pmatrix} U^*(I_p\oplus (-I_q))U & B\\ 0_k & 0_k\end{pmatrix}\Big)
\ge   w\Big(\left(\begin{array}{c|c}\begin{array}{cc} 0 & 1\\ 1 & 0\end{array} & \begin{array}{cc} s_1 & 0\\ 0 & s_k\end{array}
\\ \hline 0_2 & 0_2\end{array}\right)\oplus
\begin{pmatrix} {U'}^*(I_{p-1}\oplus (-I_{q-1}))U' & B_1\\ 0_{k-2} & 0_{k-2}\end{pmatrix}\Big) $$
for some unitary matrix $U'\in M_{k-2}$.
\end{lemma}

\begin{theorem}\label{t36}
Let $1\le q\le p$, $k=p+q$, $A=I_p\oplus(-I_q)\in M_k$, $B\in M_k$, and
$$\alpha=\min\left\{w\Big(\begin{pmatrix} U^*AU & B\\ 0_k & 0_k\end{pmatrix}\Big) : U\in M_k \mbox{ is unitary}\right\}.$$
Let  $\tilde{ A} = I_p \oplus (-I_p)\in M_{2p}$, $\tilde{B} = B \oplus s_{q+1}I_{p-q}\in M_{2p}$, and
$$\beta=\min\left\{w\Big(\begin{pmatrix} V^* \tilde{A}V & \tilde{B}\\ 0_{2p} & 0_{2p}\end{pmatrix}\Big) : V\in M_{2p} \mbox{ is unitary}\right\}.$$
Then
$$
\alpha=w\Big(\left(\begin{array}{c|c}\begin{array}{ccc}  & & J_q\\ & I_{p-q} & \\ J_q & & \end{array} & S_B
\\ \hline 0_k & 0_k\end{array}\right)\Big)
= w\Big(\begin{pmatrix}
J_{2p} & S_{\tilde{B}} \\ 0_{2p} & 0_{2p}
\end{pmatrix}\Big)=\beta,
$$
where $J_q$, $J_{2p}$, $S_B$, and $S_{\tilde{B}}$ are as in \eqref{eq1}.
\end{theorem}

\begin{proof}[Proof of Theorem \ref{t36}]
 First we consider $\alpha$.
Let $B=Q_0S_BQ_1^*$ be the singular value decomposition of $B$ and $W_0=Q_0\oplus Q_1\in M_{2k}$ be unitary, then, for any unitary $U\in M_k$, we have
$$W_0^*\begin{pmatrix} U^*AU & B\\ 0_k & 0_k\end{pmatrix}W_0=\begin{pmatrix} (UQ_0)^*A(UQ_0) & S_B\\ 0_k & 0_k\end{pmatrix}.$$
Since $UQ_0\in M_k$ is also unitary, we deduce that
\begin{equation}\label{eq3}
\alpha=\min\left\{w\Big(\begin{pmatrix} U^*AU & S_B\\ 0_k & 0_k\end{pmatrix}\Big) : U\in M_k \mbox{ is unitary}\right\}\le
w\Big(\left(\begin{array}{c|c}\begin{array}{ccc}  & & J_q\\ & I_{p-q} & \\ J_q & & \end{array} & S_B
\\ \hline 0_k & 0_k\end{array}\right)\Big).
\end{equation}
Therefore, we may assume that $B=S_B=\diag(s_1, \dots, s_k)$. Let $B_j=\diag(s_{j+1}, \dots, s_{k-j})\in M_{k-2j}$ for $j=1, \dots, q$.

Now, for any unitary $U\in M_k$, apply Lemma \ref{l35} repeatedly, we obtain
\begin{eqnarray*}
& & w\Big(\begin{pmatrix} U^*AU & B\\ 0_k & 0_k\end{pmatrix}\Big)   \\
& \ge & w\Big(\left(\begin{array}{c|c}\begin{array}{cc} 0 & 1\\ 1 & 0\end{array} & \begin{array}{cc} s_1 & 0\\ 0 & s_k\end{array}
\\ \hline 0_2 & 0_2\end{array}\right)\oplus
\begin{pmatrix} U_1^*(I_{p-1}\oplus (-I_{q-1}))U_1 & B_1\\ 0_{k-2} & 0_{k-2}\end{pmatrix}\Big) \\
& \ge & w\Big(\left(\begin{array}{c|c}\begin{array}{cc} 0 & 1\\ 1 & 0\end{array} & \begin{array}{cc} s_1 & 0\\ 0 & s_k\end{array}
\\ \hline 0_2 & 0_2\end{array}\right)\oplus
\left(\begin{array}{c|c}\begin{array}{cc} 0 & 1\\ 1 & 0\end{array} & \begin{array}{cc} s_2 & 0\\ 0 & s_{k-1}\end{array}
\\ \hline 0_2 & 0_2\end{array}\right) \oplus
\begin{pmatrix} U_2^*(I_{p-2}\oplus (-I_{q-2}))U_2 & B_2\\ 0_{k-4} & 0_{k-4}\end{pmatrix}\Big) \\
& \ge & \cdots \\
& \ge & w\Big(\left(\begin{array}{c|c}\begin{array}{cc} 0 & 1\\ 1 & 0\end{array} & \begin{array}{cc} s_1 & 0\\ 0 & s_k\end{array}
\\ \hline 0_2 & 0_2\end{array}\right)\oplus \cdots\oplus
\left(\begin{array}{c|c}\begin{array}{cc} 0 & 1\\ 1 & 0\end{array} & \begin{array}{cc} s_q & 0\\ 0 & s_{k-q+1}\end{array}
\\ \hline 0_2 & 0_2\end{array}\right) \oplus
\begin{pmatrix} U_q^*I_{p-q}U_q & B_q\\ 0_{k-2q} & 0_{k-2q}\end{pmatrix}\Big) \\
& = & w\Big(\left(\begin{array}{c|c}\begin{array}{cc} 0 & 1\\ 1 & 0\end{array} & \begin{array}{cc} s_1 & 0\\ 0 & s_k\end{array}
\\ \hline 0_2 & 0_2\end{array}\right)\oplus \cdots\oplus
\left(\begin{array}{c|c}\begin{array}{cc} 0 & 1\\ 1 & 0\end{array} & \begin{array}{cc} s_q & 0\\ 0 & s_{p+1}\end{array}
\\ \hline 0_2 & 0_2\end{array}\right) \oplus
\begin{pmatrix} I_{p-q} & B_q\\ 0_{p-q} & 0_{p-q}\end{pmatrix}\Big),
\end{eqnarray*}
where $U_j\in M_{k-2j}$ is unitary for $j=1, \dots, q$. Note that the matrix
$$\left(\begin{array}{c|c}\begin{array}{ccc}  & & J_q\\ & I_{p-q} & \\ J_q & & \end{array} & S_B
\\ \hline 0_k & 0_k\end{array}\right)$$
is permutationally similar to
$$\left(\begin{array}{c|c}\begin{array}{cc} 0 & 1\\ 1 & 0\end{array} & \begin{array}{cc} s_1 & 0\\ 0 & s_k\end{array}
\\ \hline 0_2 & 0_2\end{array}\right)\oplus \cdots\oplus
\left(\begin{array}{c|c}\begin{array}{cc} 0 & 1\\ 1 & 0\end{array} & \begin{array}{cc} s_q & 0\\ 0 & s_{p+1}\end{array}
\\ \hline 0_2 & 0_2\end{array}\right) \oplus
\begin{pmatrix} I_{p-q} & B_q\\ 0_{p-q} & 0_{p-q}\end{pmatrix},$$
hence we conclude that
$$w\Big(\begin{pmatrix} U^*AU & B\\ 0_k & 0_k\end{pmatrix}\Big) \ge w\Big(\left(\begin{array}{c|c}\begin{array}{ccc}  & & J_q\\ & I_{p-q} & \\ J_q & & \end{array} & S_B
\\ \hline 0_k & 0_k\end{array}\right)\Big)$$
for any unitary $U\in M_k$, that is,
$$\alpha\ge
w\Big(\left(\begin{array}{c|c}\begin{array}{ccc}  & & J_q\\ & I_{p-q} & \\ J_q & & \end{array} & S_B
\\ \hline 0_k & 0_k\end{array}\right)\Big) \ \ \mbox{ or } \ \
\alpha=
w\Big(\left(\begin{array}{c|c}\begin{array}{ccc}  & & J_q\\ & I_{p-q} & \\ J_q & & \end{array} & S_B
\\ \hline 0_k & 0_k\end{array}\right)\Big)$$
by \eqref{eq3}.


Now we turn to $\beta$.
From the proof on $\alpha$, we have
\begin{eqnarray*}
\alpha & = & w\Big(\left(\begin{array}{c|c}\begin{array}{ccc}  & & J_q\\ & I_{p-q} & \\ J_q & & \end{array} & S_B
\\ \hline 0_k & 0_k\end{array}\right)\Big) \\
&= &  w\Big(\left(\begin{array}{c|c}\begin{array}{cc} 0 & 1\\ 1 & 0\end{array} & \begin{array}{cc} s_1 & 0\\ 0 & s_k\end{array}
\\ \hline 0_2 & 0_2\end{array}\right)\oplus \cdots\oplus
\left(\begin{array}{c|c}\begin{array}{cc} 0 & 1\\ 1 & 0\end{array} & \begin{array}{cc} s_q & 0\\ 0 & s_{p+1}\end{array}
\\ \hline 0_2 & 0_2\end{array}\right) \oplus
\begin{pmatrix} I_{p-q} & B_q\\ 0_{p-q} & 0_{p-q}\end{pmatrix}\Big)\\
&=& \max\left\{\max_{1\le j\le q}w\Big(\left(\begin{array}{c|c}\begin{array}{cc} 0 & 1\\ 1 & 0\end{array} & \begin{array}{cc} s_j & 0\\ 0 & s_{k-j+1}\end{array}
\\ \hline 0_2 & 0_2\end{array}\right)\Big), \ \ \max_{q+1\le j\le p} w\Big(\begin{pmatrix} 1 & s_{j}\\ 0 & 0\end{pmatrix}\Big)\right\}.
\end{eqnarray*}
Since $s_{q+1}\ge s_{q+2}\ge \cdots\ge s_p$, we have
$w\Big(\begin{pmatrix} 1 & s_{j}\\ 0 & 0\end{pmatrix}\Big) \le w\Big(\begin{pmatrix} 1 & s_{q+1}\\ 0 & 0\end{pmatrix}\Big)=(1+\sqrt{1+s_{q+1}^2})/2$ for all $q+2\le j\le p$. Therefore, we obtain
\begin{equation}\label{ee1}
\alpha=\max\left\{\max_{1\le j\le q}w\Big(\left(\begin{array}{c|c}\begin{array}{cc} 0 & 1\\ 1 & 0\end{array} & \begin{array}{cc} s_j & 0\\ 0 & s_{k-j+1}\end{array}
\\ \hline 0_2 & 0_2\end{array}\right)\Big), \ \ \frac{1+\sqrt{1+s_{q+1}^2}}{2}\right\}.
\end{equation}
On the other hand, since $\tilde{B} = B \oplus s_{q+1}I_{p-q}$, from \eqref{eq1}, we have
$$S_{\tilde{B}}=\diag(s_1, \dots, s_q)\oplus s_{q+1}I_{p-q}\oplus \diag(\underbrace{s_{q+1}, \dots, s_p}_{p-q})\oplus \diag(\underbrace{s_{p+1}, \dots, s_k}_q).$$
From we have proved before, we obtain
$$
\beta  \equiv  \min\left\{w\Big(\begin{pmatrix} V^* \tilde{A}V & \tilde{B}\\ 0_{2p} & 0_{2p}\end{pmatrix}\Big) : V\in M_{2p} \mbox{ is unitary}\right\}
=  \ w\Big(\left(\begin{array}{cc} J_{2p} & S_{\tilde{B}}
\\ 0_{2p} & 0_{2p}\end{array}\right)\Big). $$
Note that $\left(\begin{array}{cc} J_{2p} & S_{\tilde{B}}
\\ 0_{2p} & 0_{2p}\end{array}\right)$ is permutationally similar to
$$\Bigg(\sum_{j=1}^q\oplus\left(\begin{array}{c|c}\begin{array}{cc} 0 & 1\\ 1 & 0\end{array} & \begin{array}{cc} s_j & 0\\ 0 & s_{k-j+1}\end{array}
\\ \hline 0_2 & 0_2\end{array}\right)\Bigg) \oplus \Bigg(\sum_{j=q+1}^p\oplus\left(\begin{array}{c|c}\begin{array}{cc} 0 & 1\\ 1 & 0\end{array} & \begin{array}{cc} s_{q+1} & 0\\ 0 & s_{j}\end{array}
\\ \hline 0_2 & 0_2\end{array}\right)\Bigg).$$
In particular, the matrix
$$\left(\begin{array}{c|c}\begin{array}{cc} 0 & 1\\ 1 & 0\end{array} & \begin{array}{cc} s_{q+1} & 0\\ 0 & s_{q+1}\end{array}
\\ \hline 0_2 & 0_2\end{array}\right)$$ is unitarily similar to $\begin{pmatrix} 1 & s_{q+1} \\ 0 & 0\end{pmatrix}\oplus \begin{pmatrix} -1 & s_{q+1} \\ 0 & 0\end{pmatrix}$, this yields that
$$w\Big(\left(\begin{array}{c|c}\begin{array}{cc} 0 & 1\\ 1 & 0\end{array} & \begin{array}{cc} s_{q+1} & 0\\ 0 & s_{q+1}\end{array}
\\ \hline 0_2 & 0_2\end{array}\right)\Big)
=w\Big( \begin{pmatrix} 1 & s_{q+1} \\ 0 & 0\end{pmatrix}\Big)
=  \frac{1+\sqrt{1+s_{q+1}^2}}{2}.$$
Moreover, since $s_{q+1}\ge s_{q+2}\ge \cdots\ge s_p$, by Corollary \ref{2.4}, we have
$$w\Big(\left(\begin{array}{c|c}\begin{array}{cc} 0 & 1\\ 1 & 0\end{array} & \begin{array}{cc} s_{q+1} & 0\\ 0 & s_{j}\end{array}
\\ \hline 0_2 & 0_2\end{array}\right)\Big) \le w\Big(\left(\begin{array}{c|c}\begin{array}{cc} 0 & 1\\ 1 & 0\end{array} & \begin{array}{cc} s_{q+1} & 0\\ 0 & s_{q+1}\end{array}
\\ \hline 0_2 & 0_2\end{array}\right)\Big) \ \  \mbox{for all }  q+2\le j\le p.$$
It follows that
$$w\Big(\sum_{j=q+1}^p\oplus\left(\begin{array}{c|c}\begin{array}{cc} 0 & 1\\ 1 & 0\end{array} & \begin{array}{cc} s_{q+1} & 0\\ 0 & s_{j}\end{array}
\\ \hline 0_2 & 0_2\end{array}\right)\Big)= w\Big(\left(\begin{array}{c|c}\begin{array}{cc} 0 & 1\\ 1 & 0\end{array} & \begin{array}{cc} s_{q+1} & 0\\ 0 & s_{q+1}\end{array}
\\ \hline 0_2 & 0_2\end{array}\right)\Big) = \frac{1+\sqrt{1+s_{q+1}^2}}{2}.$$
Hence we obtain
\begin{eqnarray*}
\beta &= & w\Big(\left(\begin{array}{cc} J_{2p} & S_{\tilde{B}}
\\ 0_{2p} & 0_{2p}\end{array}\right)\Big) \\
& = & \max\left\{ w\Big(\sum_{j=1}^q\oplus\left(\begin{array}{c|c}\begin{array}{cc} 0 & 1\\ 1 & 0\end{array} & \begin{array}{cc} s_j & 0\\ 0 & s_{k-j+1}\end{array}
\\ \hline 0_2 & 0_2\end{array}\right)\Big), w\Big( \sum_{j=q+1}^p\oplus\left(\begin{array}{c|c}\begin{array}{cc} 0 & 1\\ 1 & 0\end{array} & \begin{array}{cc} s_{q+1} & 0\\ 0 & s_{j}\end{array}
\\ \hline 0_2 & 0_2\end{array}\right)\Big)  \right\}\\
& = & \max\left\{\max_{1\le j\le q}w\Big(\left(\begin{array}{c|c}\begin{array}{cc} 0 & 1\\ 1 & 0\end{array} & \begin{array}{cc} s_j & 0\\ 0 & s_{k-j+1}\end{array}
\\ \hline 0_2 & 0_2\end{array}\right)\Big), \ \ \frac{1+\sqrt{1+s_{q+1}^2}}{2}\right\} = \alpha
\end{eqnarray*}
from \eqref{ee1}, this completes the proof.
\end{proof}

\end{document}
